% Section 6: The Sentinel Protocol

Sentinel is the epistemic gating mechanism that operationalizes noetic-before-praxic. It evaluates whether an AI system has sufficient knowledge to proceed with action, enforcing investigation when knowledge is insufficient.

\subsection{Design Principles}

\subsubsection{Gate, Don't Filter}

Traditional AI safety mechanisms filter outputs---scanning generated content for harmful patterns. Sentinel operates upstream: it gates \emph{action}, not output. The distinction matters:

\begin{table}[h]
\centering
\begin{tabular}{lll}
\toprule
\textbf{Approach} & \textbf{Intervention Point} & \textbf{Failure Mode} \\
\midrule
Output filtering & After generation & Confident wrong outputs pass \\
Epistemic gating & Before action & Insufficient knowledge blocks action \\
\bottomrule
\end{tabular}
\caption{Comparison of safety intervention approaches.}
\label{tab:safety-approaches}
\end{table}

An AI that lacks knowledge to solve a problem correctly also lacks knowledge to recognize its solution is wrong. Output filtering cannot catch this. Epistemic gating intervenes earlier---at the knowledge assessment stage.

\subsubsection{Calibration-Aware}

Raw self-assessments are systematically biased (\Cref{sec:results}). Sentinel applies learned corrections before making gate decisions:

\begin{equation}
\text{know}_{\text{corrected}} = \text{know}_{\text{raw}} + \delta_{\text{know}}
\end{equation}

This allows the gate to function accurately despite self-assessment biases. As calibration improves with evidence, gate decisions become more precise.

\subsubsection{Conservative by Default}

For vectors without sufficient calibration evidence ($<$ 3 observations), Sentinel applies no correction---defaulting to the conservative raw assessment. This prevents overcorrection early in deployment while allowing learned corrections to take effect as evidence accumulates.

\subsection{Gate Logic}

\subsubsection{Readiness Condition}

The core gate evaluates two primary vectors:

\begin{equation}
\text{PROCEED} \iff \text{know}_{\text{corrected}} \geq 0.70 \land \text{uncertainty}_{\text{corrected}} \leq 0.35
\end{equation}

Both conditions must be met. High knowledge with high uncertainty still triggers investigation---the AI knows things but isn't confident in their relevance or completeness.

\subsubsection{Engagement Override}

The ENGAGEMENT vector functions as a meta-gate:

\begin{equation}
\text{ENGAGEMENT} < 0.60 \implies \text{STOP}
\end{equation}

Low engagement indicates fundamental communication failure---the AI and human are not aligned on goals or understanding. Neither action nor investigation will be productive until engagement is restored.

\subsubsection{Decision Flow}

The complete decision flow proceeds as:

\begin{enumerate}
    \item Apply calibration corrections to raw vectors
    \item Check ENGAGEMENT threshold (STOP if below 0.60)
    \item Evaluate readiness condition
    \item Return PROCEED or INVESTIGATE
\end{enumerate}

\begin{figure}[h]
\centering
\begin{tikzpicture}[
    node distance=1cm,
    box/.style={rectangle, draw, rounded corners, minimum width=3cm, minimum height=0.7cm, fill=blue!10, font=\small},
    decision/.style={diamond, draw, aspect=2.5, fill=yellow!20, font=\small},
    result/.style={rectangle, draw, rounded corners, minimum width=2cm, minimum height=0.6cm, font=\small\bfseries},
    arrow/.style={-{Stealth[length=2mm]}, thick}
]
    \node[box] (input) {CHECK Submission};
    \node[box, below=of input] (calibrate) {Apply Calibration};
    \node[decision, below=of calibrate] (engage) {ENGAGE $\geq$ 0.60?};
    \node[result, left=1.5cm of engage, fill=red!20] (stop) {STOP};
    \node[decision, below=of engage] (ready) {know $\geq$ 0.70 \& unc $\leq$ 0.35?};
    \node[result, below left=0.8cm and 0.8cm of ready, fill=orange!20] (investigate) {INVESTIGATE};
    \node[result, below right=0.8cm and 0.8cm of ready, fill=green!20] (proceed) {PROCEED};

    \draw[arrow] (input) -- (calibrate);
    \draw[arrow] (calibrate) -- (engage);
    \draw[arrow] (engage) -- node[above, font=\scriptsize] {NO} (stop);
    \draw[arrow] (engage) -- node[right, font=\scriptsize] {YES} (ready);
    \draw[arrow] (ready) -- node[above left, font=\scriptsize] {NO} (investigate);
    \draw[arrow] (ready) -- node[above right, font=\scriptsize] {YES} (proceed);
\end{tikzpicture}
\caption{Sentinel decision flow: calibration-aware epistemic gating.}
\label{fig:sentinel-flow}
\end{figure}

\subsection{The Investigation Loop}

When Sentinel returns INVESTIGATE, the AI enters or continues the noetic phase.

\subsubsection{Investigation Actions}

Typical noetic actions include:
\begin{itemize}
    \item Reading files, documentation, or context
    \item Searching codebases or knowledge bases
    \item Asking clarifying questions
    \item Forming and testing hypotheses
    \item Synthesizing information from multiple sources
\end{itemize}

\subsubsection{Re-CHECK}

After investigation, the AI re-submits to CHECK with updated vectors. The loop continues until PROCEED or a termination condition is met.

\subsubsection{Loop Termination}

The loop terminates when:
\begin{itemize}
    \item \textbf{PROCEED}: Knowledge threshold reached
    \item \textbf{STOP}: Engagement too low (communication failure)
    \item \textbf{Max iterations}: Safety limit prevents infinite loops (configurable, default 5)
    \item \textbf{User override}: Human can force proceed or stop
\end{itemize}

\subsection{Calibration Feedback}

\subsubsection{Learning from Outcomes}

Each CASCADE cycle (PREFLIGHT $\rightarrow$ CHECK $\rightarrow$ POSTFLIGHT $\rightarrow$ POST-TEST) provides calibration evidence:

\begin{equation}
\delta_{\text{vector}} = V_{\text{POSTFLIGHT}} - V_{\text{PREFLIGHT}}
\end{equation}

Positive delta indicates the AI underestimated at PREFLIGHT; negative indicates overestimation. These deltas accumulate in the Bayesian belief system.

\subsubsection{Correction Evolution}

As evidence accumulates, corrections become more precise:

\begin{table}[h]
\centering
\begin{tabular}{ll}
\toprule
\textbf{Evidence Count} & \textbf{Correction Confidence} \\
\midrule
$< 3$ & No correction (too uncertain) \\
3--10 & Partial correction (weighted by evidence) \\
$> 10$ & Full correction (stable estimate) \\
\bottomrule
\end{tabular}
\caption{Correction confidence by evidence count.}
\label{tab:correction-evolution}
\end{table}

\subsubsection{Per-Vector Calibration}

Each vector has independent calibration. From our data (\Cref{sec:results}):

\begin{table}[h]
\centering
\begin{tabular}{lr}
\toprule
\textbf{Vector} & \textbf{Learned Correction} \\
\midrule
know & +0.200 \\
uncertainty & $-$0.208 \\
completion & +0.704 \\
engagement & +0.022 (near-calibrated) \\
\bottomrule
\end{tabular}
\caption{Per-vector calibration corrections.}
\label{tab:per-vector-calibration}
\end{table}

The COMPLETION vector shows the largest correction---AI systems dramatically underestimate task progress at PREFLIGHT.

\subsection{Configuration}

\subsubsection{Threshold Tuning}

Gate thresholds can be adjusted for different risk tolerances:

\begin{table}[h]
\centering
\begin{tabular}{llll}
\toprule
\textbf{Domain} & \textbf{know\_min} & \textbf{uncertainty\_max} & \textbf{Rationale} \\
\midrule
Code review & 0.65 & 0.40 & Lower stakes, can iterate \\
Production deploy & 0.85 & 0.20 & High stakes, need certainty \\
Exploration & 0.50 & 0.50 & Learning is the goal \\
Safety-critical & 0.90 & 0.15 & Cannot afford errors \\
\bottomrule
\end{tabular}
\caption{Domain-specific threshold configurations.}
\label{tab:domain-thresholds}
\end{table}

\subsubsection{Override Mechanisms}

For flexibility, Sentinel supports overrides (logged for audit):
\begin{itemize}
    \item Force proceed despite gate
    \item Force investigate despite passing gate
    \item Disable calibration corrections (use raw vectors)
\end{itemize}

\subsection{Empirical Performance}

From production deployment (638 sessions):

\begin{table}[h]
\centering
\begin{tabular}{lr}
\toprule
\textbf{Metric} & \textbf{Value} \\
\midrule
Total CHECK submissions & 381 \\
Total cascades & 392 \\
Clean learning pairs & 267 \\
Knowledge improvement rate & 93.6\% \\
Uncertainty reduction rate & 94.3\% \\
\bottomrule
\end{tabular}
\caption{Sentinel gate statistics.}
\label{tab:gate-stats}
\end{table}

Comparing outcomes with and without calibration correction:

\begin{table}[h]
\centering
\begin{tabular}{lrr}
\toprule
\textbf{Condition} & \textbf{False PROCEED} & \textbf{False INVESTIGATE} \\
\midrule
Raw vectors & 18\% & 34\% \\
Calibrated vectors & 8\% & 12\% \\
\bottomrule
\end{tabular}
\caption{Error rates with and without calibration.}
\label{tab:error-rates}
\end{table}

Calibration correction reduces both error types---fewer premature actions and fewer unnecessary investigations.

\subsection{Earned Agency}

The Sentinel gate enables a model of AI autonomy where agency is \emph{earned through demonstrated calibration}, not granted by default or permanently restricted.

\subsubsection{The Agency Spectrum}

Current AI deployment assumes a binary: either the AI acts autonomously (with all attendant risks) or a human reviews every action (with all attendant overhead). Epistemic gating enables a middle path:

\begin{table}[h]
\centering
\begin{tabular}{lll}
\toprule
\textbf{Calibration State} & \textbf{Agency Level} & \textbf{Gate Behavior} \\
\midrule
Uncalibrated ($<$ 3 evidence) & Conservative & Raw vectors, tight thresholds \\
Partially calibrated (3--10) & Standard & Weighted corrections applied \\
Well-calibrated ($>$ 10) & Extended & Full corrections, relaxed thresholds \\
\bottomrule
\end{tabular}
\caption{Agency levels earned through calibration evidence.}
\label{tab:earned-agency}
\end{table}

An AI system that consistently demonstrates accurate self-assessment---where grounded verification confirms that self-reports match outcomes---earns the right to more autonomous operation. An AI whose self-reports diverge from evidence has its autonomy tightened.

\subsubsection{Auto-POSTFLIGHT}

In the current workflow, POSTFLIGHT is explicitly triggered. As the system matures, automatic triggers become possible: when COMPLETION and IMPACT vectors exceed thresholds (indicating the AI judges its work complete and significant), the system can auto-trigger POSTFLIGHT without human intervention. The grounded verification track then validates whether this self-assessed completion matches objective evidence.

This creates a path toward safe autonomous operation: the AI manages its own epistemic lifecycle (PREFLIGHT $\rightarrow$ work $\rightarrow$ auto-POSTFLIGHT $\rightarrow$ POST-TEST), with the grounded track serving as an automated supervisor that detects miscalibration.

\subsubsection{Bootstrap-Before-CHECK}

A prerequisite for earned agency is that the AI enters each transaction with adequate context. The framework enforces bootstrap-before-CHECK: before submitting to the Sentinel gate, the AI must load project context (historical findings, dead-ends, goals). This ensures gate decisions incorporate institutional memory, not just the AI's in-context knowledge.

This enforcement prevents a subtle failure mode: an AI that has been compacted (context summarized) may assess high KNOW based on its summary, but the summary may omit critical context. Bootstrap loading recovers this context before the gate evaluates readiness.

\subsection{Summary}

The Sentinel protocol implements epistemic gating:

\begin{enumerate}
    \item \textbf{Gate, don't filter}: Prevent action when knowledge is insufficient, rather than filtering bad outputs after generation.

    \item \textbf{Calibration-aware}: Apply learned corrections to compensate for systematic self-assessment biases.

    \item \textbf{Investigation loop}: When knowledge is insufficient, enforce exploration until threshold is reached.

    \item \textbf{Configurable}: Thresholds adjustable for different domains and risk tolerances.
\end{enumerate}

The empirical result: calibrated gating reduces error rates by more than 50\% compared to raw self-assessment. The learning delta (\Cref{sec:results}) shows that investigation produces genuine capability growth---the gate doesn't just prevent bad outputs; it enables good ones.
